The complexity class at the bottom of our hierarchy is the class \textbf{P}. The class \textbf{P} is defined as follows:
$$ P = \bigcup_{c \geq 1} \textbf{DTIME}(n^c)$$
We get the union of all the \textbf{DTIME}$(n^c)$ classes whose function is any polynomial time representation $n^c$, therefore obtaining an unique set of all 
those tasks sharing the same time complexity.

In other words, the class \textbf{P} includes all those languages $L$ (decision problems) that can be decided (solved) by a Turing Machine in a running time $P$,
where $P$ can be any polynomial time representation of the form $n^c$.

More formally, for any polynomial \textbf{P} there are $c,d$ natural numbers that are strictly greater than 0 such that $P(n) \leq c \cdot n^d$ for suffiently 
large $n$. By this last statement, we can observe that $c$ and $d$ can be arbitraly large or huge, so, for instance, a \textit{Turing Machine} deciding a given 
language $L$ (decision problem) in time $10^{20} \cdot n^{10^{30}}$ still always a resolution in a polynomial time!

\textbf{P} is generally considered as \textbf{the class} of efficiently decidable languages, problems or tasks of this class seem to be simple, but anyway
it's always possible to build a machine that solves the same problem or task with a different complexity.